\section{Производная}

\begin{frame}[c]{Таблица производных и правила}
  \begin{columns}[T]
    \begin{column}{0.40\textwidth}
      \centering
      \zebra{2}
      \begin{tabular}{cc}
        \textbf{$f(x)$} & \textbf{$f'(x)$}\\
        $c$ & $0$\\
        $x^n$ & $n\,x^{n-1}$\\
        $e^x$ & $e^x$\\
        $\ln x$ & $\dfrac1x$\\
        $\sin x$ & $\cos x$\\
        $\cos x$ & $-\sin x$\\
      \end{tabular}
    \end{column}
    \begin{column}{0.54\textwidth}
      \begin{block}{Правила}
        \small
        \begin{tabular}{@{}ll@{}}
          сумма & $(f+g)' = f' + g'$\\[3pt]
          число & $(c\,f)' = c\,f'$\\[3pt]
          произведение & $(f g)' = f' g + f g'$\\[3pt]
          \multicolumn{2}{@{}l}{\textbf{сложная функция}}\\[2pt]
          \multicolumn{2}{@{}l}{\quad$\bigl(f(g(x))\bigr)' = f'\bigl(g(x)\bigr)\cdot g'(x)$}\\
        \end{tabular}
      \end{block}
      \begin{block}{Цепное правило словами}
        \small
        Производная внешней функции, внутренность не трогаем, умножить на производную
        внутренней.
      \end{block}
    \end{column}
  \end{columns}
\end{frame}

\begin{frame}[c]{Разминка: правила}
  \lead{Самостоятельно.} Найдите производные.

  \medskip
  \begin{numrows}
    1 & $f(x) = 4x^3 - 2x + 7$\\
    2 & $f(x) = 3\sin x + e^x$\\
  \end{numrows}
\end{frame}

\begin{frame}[c]{Ответы: разминка}
  \begin{numrows}
    1 & $12x^2 - 2$\\
    2 & $3\cos x + e^x$\\
  \end{numrows}
\end{frame}

\begin{frame}[c]{Цепное правило на примерах}
  \zebra{2}
  \begin{tabular}{lll}
    \textbf{Функция} & \textbf{Внешняя и внутренняя} & \textbf{Производная}\\
    $(3x-1)^2$ & $u^2$, \ $u = 3x-1$ & $2(3x-1)\cdot 3 = 6(3x-1)$\\
    $\sin^2 x$ & $u^2$, \ $u = \sin x$ & $2\sin x\cos x$\\
    $\ln(x^2+1)$ & $\ln u$, \ $u = x^2+1$ & $\dfrac{2x}{x^2+1}$\\
    $e^{-x^2}$ & $e^u$, \ $u = -x^2$ & $-2x\,e^{-x^2}$\\
  \end{tabular}

  \medskip
  {\footnotesize\color{acGray}Проверить себя можно численно:
  $f'(x) \approx \frac{f(x+h)-f(x-h)}{2h}$ при маленьком $h$, например $10^{-5}$.\par}
\end{frame}

\begin{frame}[c]{Самостоятельно: цепное правило}
  \begin{numrows}
    1 & $f(x) = (2x+5)^3$\\
    2 & $f(x) = x\,e^{2x}$ \quad{\footnotesize\color{acGray}(произведение и цепное правило)}\\
    $\bigstar$ & $\sigma(x) = \dfrac{1}{1 + e^{-x}}$. Покажите, что $\sigma'(x) = \sigma(x)\bigl(1 - \sigma(x)\bigr)$\\
  \end{numrows}
\end{frame}

\begin{frame}[c]{Ответы: цепное правило}
  \begin{numrows}
    1 & $3(2x+5)^2\cdot 2 = 6(2x+5)^2$\\
    2 & $e^{2x} + 2x\,e^{2x} = (1 + 2x)\,e^{2x}$\\
    $\bigstar$ & $\sigma' = \dfrac{e^{-x}}{(1+e^{-x})^2} = \dfrac{1}{1+e^{-x}}\cdot\dfrac{e^{-x}}{1+e^{-x}} = \sigma\,(1-\sigma)$\\
  \end{numrows}
\end{frame}

\begin{frame}[c]{Производная по параметру, а не по $x$}
  \lead{В ML переменная --- это параметр модели, а данные $x$, $y$ --- просто числа.}

  \medskip
  \begin{columns}[T]
    \begin{column}{0.47\textwidth}
      \begin{block}{Ошибка на одном объекте}
        \small
        \[\frac{\partial}{\partial w}(w x - y)^2 = 2(wx-y)\cdot x\]
        Внешняя функция --- квадрат, внутренняя --- $wx-y$. Её производная по $w$ равна $x$.
      \end{block}
    \end{column}
    \begin{column}{0.47\textwidth}
      \begin{block}{Два признака}
        \small
        \[\frac{\partial}{\partial w_2}(w_1 x_1 + w_2 x_2 - y)^2 = 2(w_1 x_1 + w_2 x_2 - y)\cdot x_2\]
        $w_1$ при этом считаем числом: производная по $w_2$ от $w_1 x_1$ равна нулю.
      \end{block}
    \end{column}
  \end{columns}
\end{frame}

\begin{frame}[c]{Самостоятельно: производная по параметру}
  \begin{numrows}
    1 & $L(b) = (3w + b - 7)^2$ по $b$\\
    2 & $L(w) = e^{w x}$ по $w$\\
  \end{numrows}
\end{frame}

\begin{frame}[c]{Ответы: производная по параметру}
  \begin{numrows}
    1 & $2(3w + b - 7)$\\
    2 & $x\,e^{wx}$: здесь $x$ --- число\\
  \end{numrows}
\end{frame}
